%% ma: L12-B12-0003
%% lop: 12
%% chuong: 4
%% bai: 12
%% dang: 12.12.2
%% loai: DS
%% muc: [TH, VD, VD, VD]
%% dap_an: "ĐĐĐS"
%% nguon: "(NBV)-ĐỀ SỐ 1-MỖI NGÀY 1 ĐỀ THI 2025.pdf (D01, MỖI NGÀY 1 ĐỀ - Đề số 1), Phần II Câu 2 (THCS THPT Lê Thánh Tông - HCM 2025)"
%% kien_thuc: "Bài 12 (tích phân, quãng đường từ vận tốc), Bài 11 (nguyên hàm); bài toán chuyển động chậm dần đều (vận tốc là nguyên hàm của gia tốc)"
%% ghi_chu: "Trùng ý tưởng với câu cùng dạng, giáo viên đã duyệt giữ (09/10/2026). Đề gốc có ảnh minh họa đường cao tốc (không cần cho lời giải). Đáp án gốc Đ Đ Đ S đúng (kiểm bằng sympy)"
%% trang_thai: chinh-thuc
\begin{ex}
Một người điều khiển ô tô đang ở đường dẫn muốn nhập làn vào đường cao tốc. Khi ô tô cách điểm nhập làn $200$ m, tốc độ của ô tô là $36$ km/h. Hai giây sau đó, ô tô bắt đầu tăng tốc với tốc độ $v(t)=at+b$ ($a,b\in\mathbb R$, $a>0$), trong đó $t$ là thời gian tính bằng giây kể từ khi bắt đầu tăng tốc. Biết rằng ô tô nhập làn cao tốc sau $12$ giây tăng tốc và duy trì sự tăng tốc trong $24$ giây kể từ khi bắt đầu tăng tốc, sau $24$ s đó ô tô duy trì tốc độ cao nhất trong $24$ giây đầu trong thời gian còn lại trên cao tốc.
\choiceTF
{\True Quãng đường ô tô đi được từ khi bắt đầu tăng tốc đến khi nhập làn là $180$ m.}
{\True Vận tốc của ô tô tại thời điểm nhập làn là $72$ km/h.}
{\True Quãng đường mà ô tô đi được trong thời gian $30$ giây kể từ khi ô tô cách điểm nhập làn $200$ m là $620$ m.}
{Sau $24$ giây kể từ khi tăng tốc, ô tô duy trì tốc độ cao nhất trong vòng $5$ s thì phát hiện chướng ngại vật cách đó $300$ m, người điều khiển lập tức đạp phanh và ô tô chuyển động chậm dần đều với $a(t)=-3\ (\mathrm{m/s^2})$. Khi đó ô tô dừng lại cách chướng ngại vật $10$ m.}
\loigiai{$36$ km/h $=10$ m/s nên $b=10$. Trong $2$ giây đầu ô tô đi được $20$ m, còn lại $180$ m đến điểm nhập làn.
a) Đúng. b) $\displaystyle\int_0^{12}(at+10)\,\mathrm dt=72a+120=180\Rightarrow a=\dfrac56$; $v(12)=\dfrac56\cdot12+10=20$ m/s $=72$ km/h. Đúng.
c) $v(24)=\dfrac56\cdot24+10=30$ m/s. Trong $30$ s: $2$ s đầu đi $20$ m; $24$ s tăng tốc đi $\displaystyle\int_0^{24}\left(\dfrac56t+10\right)\mathrm dt=480$ m; $4$ s còn lại đi với $30$ m/s được $120$ m. Tổng $20+480+120=620$ m. Đúng.
d) Lúc phanh $v_0=30$ m/s; $v(t)=30-3t$ dừng sau $10$ s, quãng đường phanh $\displaystyle\int_0^{10}(30-3t)\,\mathrm dt=150$ m. Ô tô dừng cách chướng ngại vật $300-150=150$ m $\ne10$ m. Sai.}
\end{ex}
